\section{Ground undecidability within additive distance one}
\label{sec:control-hardness}

The preceding constructions require growth and literal source identities.
They do not require any periodicity of the control. This permits a sharp
contrast with the ground-decision controls of \cite{groundcomparison}.

Fix an effective enumeration of programs and a primitive recursive
bounded simulation predicate $H(e,t)$, equal to one exactly when the
blank-input program with index $e$ has halted within $t$ steps.
Blank-input halting is undecidable: a decision procedure would also
decide whether a program halts on its own index, by hardwiring that
input. Define a program $D(e)$ which halts exactly when this procedure
says that program $e$ does not halt on input $e$. Evaluating $D$ at
its own index gives the diagonal contradiction. The
computability background originates with Turing \cite{turing}.

For a fixed $b\in\N$ define
\[
 \epsilon_b(0)=0,\qquad
 \epsilon_b(n)=\begin{cases}
 0,&0<n\le b,\\ H(v_2(n),n),&n>b,
 \end{cases}
 \qquad \Gamma_b(n)=G(n)+\epsilon_b(n),
\]
where $v_2(n)$ is the exponent of two dividing $n>0$.
The single function $\Gamma_b$ is fixed before a program index is
supplied to a comparison query.

\begin{lemma}[Permitted growth and a bounded perturbation]
\label{lem:perturbed-control}
The control $\Gamma_b$ is primitive recursive and agrees with $G$
on $[0,b]$. It satisfies
\[
 0\le\Gamma_b(n)-G(n)\le1,\qquad
 G(n)\le\Gamma_b(n)<G(n+1),\qquad
 \Gamma_b(n+1)\ge2^{2^{\Gamma_b(n)}}.
\]
It is eventually bounded by $F_6$, has the additive margin
$\Gamma_b(v+a)\ge\Gamma_b(v)+2a$ for $a\ge1$, and
$\Gamma_b(n)/G(n)\longrightarrow1$.
\end{lemma}
\begin{proof}
Bounded simulation, $v_2$ on positive integers and each fixed $F_j$
are primitive recursive. The prefix and additive-distance assertions
are immediate. Put $u=G(n)$ and $s=n+100$. Strict increase of the
iterates gives
\[
 G(n+1)=F_4^{s+2}(s+1)
       \ge F_4(G(n)+1)\ge2^{2^{u+1}}.
\]
The last bound follows from $F_4(x)\ge F_1(x)$ and
$F_1(x)=F_0^{x+1}(x)\ge F_0^2(x)$ at these positive inputs.
Since $\Gamma_b(n)\le u+1$ and $\Gamma_b(n+1)\ge G(n+1)$,
this proves the required adjacent estimate. It also proves
$G(n)+1<G(n+1)$ and an increment greater than two; summing those
increments gives the additive margin.

For $n\ge101$, $F_5(n)\ge n+101$, so
\[
 \Gamma_b(n)<F_5(n+101)\le F_5(F_5(n))\le F_6(n).
\]
Finally the relative error is between zero and $1/G(n)$, which
tends to zero.
\end{proof}

In particular the entire source interface and the eraser construction
apply at $f=\Gamma_b$. Its successor graph is
$F_S(z)=\Gamma_b(z+4)$. The strong growth estimate also gives the
fourfold separation of $g(n)=\Gamma_b(n+2)+1$, including its added
unit. No residue theorem for $G$ is transferred to $\Gamma_b$.

\begin{theorem}[Ground hardness at one fixed control]
\label{thm:control-hardness}
For every fixed $b$, universal canonical weak and strict comparison at
$\Gamma_b$ are $\Pi^0_1$-complete already for ground-output terms
with only ground free variables and ground binders, prime constants
$0,S$, and internal type levels at most one.
Moreover $\Gamma_b(n)\bmod2$ is not eventually periodic.
\end{theorem}
\begin{proof}
Every $G(n)$ is a positive power of two and is even, by its
definition through positive iterations of $F_0$. Thus, for $n>b$,
$\Gamma_b(n)$ is odd exactly when $H(v_2(n),n)=1$.

Fix the odd integer $C=2b+5$. For a program index $e$ form the
existential conjunction over natural variables
\[
 z+4=2^e(2t+C),\qquad y=F_S(z),\qquad y=w+w+1.
\]
The coefficient $2^e$ is a fixed integer computed from the index;
there is no variable exponentiation operation. Set
$n=2^e(2t+C)$. Then $n\ge5$, $n>b$ and $v_2(n)=e$.
At each $t$ the conjunction is solvable exactly when $H(e,n)=1$.
If the program halts, choose $t$ so that $n$ exceeds the halting
time. Then $z=n-4\ge0$ and $w=(\Gamma_b(n)-1)/2\ge0$
give witnesses. Conversely every solution certifies a halting event.

Flatten using the source addition, constant and actual $S$ graph
gadgets, and apply the collision detector and conjunction construction.
Unsatisfiability is one universal strict comparison and, by $KR\mle IL$,
one weak comparison. All newly introduced binders and variables are
ground, and the only prime constants are $0,S$. Their types have
level at most one. This is a many-one reduction of nonhalting. Its
arithmetic graph has $O(e+\log(b+2))$ nodes, a bound in the numerical
index $e$, not in its binary length.

For membership, enumerate finite tuples of positive raw ground
values. Each source value is computable. Indeed in this restricted
grammar its positive fields are finite, and each scalar collapse is
a fixed finite iterate of the primitive recursive control, by the
finite-coefficient formula. A failed universal comparison is therefore
semidecidable. Nonhalting is $\Pi^0_1$-complete: for a computable
predicate $R(e,t)$, a program can search for the first failure of
$R(e,t)$, so its nonhalting expresses $\forall t\,R(e,t)$.
This proves the stated completeness.

Suppose, finally, that parity had period $p\ge1$ after $N$.
For each $e$ effectively choose $n_e>\max(N,b)$ with $v_2(n_e)=e$.
All integers $n_e+kp2^{e+1}$ have the same residue modulo $p$ and
the same valuation $e$. The bounded tests
$H(e,n_e+kp2^{e+1})$ would therefore be constant. They are always
zero when $e$ never halts and eventually one when it halts.
Computing $\Gamma_b(n_e)\bmod2$ would decide halting. The alleged
finite $N,p$ may be hardwired; a uniform procedure finding them is
not required for the contradiction. No eventual parity period exists.
\end{proof}

The theorem changes the control and keeps the source fragment small.
It is distinct from functional-input hardness at the unchanged standard
control $G$, proved in \cite{groundcomparison}. Neither theorem settles the
all-ground problem for the exact standard function $G$.
